% RESULTADO_RECUPERADO. II REV10, 07_elipse.tex:1--61,
% 07b_geometria_elipse.tex:73--115 y 12_dualidad_t.tex:27--107.
% Explicitaciones algebraicas: cota H5<2 para verificar la carta y
% sustitución Rstar^4=(499 alpha-eta_ret)/(501 alpha+eta_ret).
\section{Ellipse d’action et détermination du rayon normalisé}
\label{iv:ellipse:section}

Le lecteur angulaire et la section d’action conservent le même
antécédent APP--TRIT--TPK. La forme construite
ci-après utilise les représentations
\eqref{iv:angular:representantes} et une section positive
\(\hbar=\hbar_s\), avec
\(s\in\{\mathrm{pre},\mathrm{ret}\}\), de
\eqref{iv:action:hbar}. Le rapport angulaire déterminera la forme ;
la section d’action conservera son échelle commune.

\subsection{Domaine et demi-axes étiquetés
}
\label{iv:ellipse:dominio}

Pour abréger, soient \(A=A_{\deg}\) et \(C^*=C^*_{\deg}\).
Dans le système de coordonnées \(A\ne0\), \(\hbar>0\) et \(|C^*/A|<1\),
nous définissons
\begin{equation}
 \xi=\frac{C^*}{A},\qquad
 \eta_{\mathrm{el}}=\operatorname{artanh}\xi,\qquad
 a_S=\sqrt{2\hbar}\,e^{\eta_{\mathrm{el}}/2},\qquad
 b_S=\sqrt{2\hbar}\,e^{-\eta_{\mathrm{el}}/2}.
 \label{iv:ellipse:semiejes}
\end{equation}
Les demi-axes et les coordonnées équilibrées \(Q,P\) ont pour
unité la racine d’action. Les axes restent étiquetés :
le signe de \(\eta_{\mathrm{el}}\) détermine lequel est le
plus grand. La courbe et son produit de demi-axes sont

\begin{equation}
 \gamma(\theta)=(a_S\cos\theta,b_S\sin\theta),\qquad
 \frac{Q^2}{a_S^2}+\frac{P^2}{b_S^2}=1,\qquad
 a_Sb_S=2\hbar.
 \label{iv:ellipse:curva}
\end{equation}

\begin{proposition}[Appartenance des représentations au système de coordonnées elliptique]
\label{iv:ellipse:pertenencia}
La branche positive de \eqref{iv:angular:representantes}, avec les
intervalles \eqref{iv:action:intervalos}, vérifie
\(A>0\), \(0<C^*<A\) et \(\hbar_s>0\).
Le renversement de la coordonnée impaire conserve \(|C^*/A|<1\).
\end{proposition}
\begin{proof}
La proposition \ref{iv:action:positividad} donne
\(0<\eta_{\mathrm{ret}}<H_5\) et la positivité des deux
sections d’action. Pour une borne supérieure de \(H_5\),
\eqref{iv:action:intervalos} implique
\(\alpha<1/125\) et \(\varphi>8/5\). En conservant seulement
les termes positifs de \eqref{iv:action:h5}, on obtient
\[
 H_5<
 \frac{\sqrt5}{2}+
 \frac{9}{16\cdot125^2}+
 \frac{7}{48\cdot125^4}.
\]
Comme \(5<81/16\), on a \(\sqrt5/2<9/8\). De plus,

\[
 \frac{9}{16\cdot125^2}+
 \frac{7}{48\cdot125^4}
 =\frac{421882}{11718750000}<\frac1{1000}.
\]
Par conséquent, \(H_5<9/8+1/1000<2\). En utilisant maintenant
\(\alpha>7/1000\),
\[
 0<C^*=2(\eta_{\mathrm{ret}}+\alpha)
 <2\left(2+\frac1{125}\right)=\frac{502}{125}
 <7<1000\alpha=A.
\]
L’échange de canaux de
\ref{iv:angular:inversion} remplace \(C^*\) par \(-C^*\)
et conserve \(A\) ; il conserve donc l’inégalité absolue.

\end{proof}

\subsection{Phase paramétrique et aire orientée}
\label{iv:ellipse:area}

\begin{proposition}[Loi aréale et transport des coordonnées angulaires]
\label{iv:ellipse:ley-areal}
La phase paramétrique \(\theta\) de \(\gamma\) et son angle
polaire \(\phi\), relevés continûment à partir de zéro, satisfont
\begin{equation}
 \phi(\theta)=\arg(a_S\cos\theta+i b_S\sin\theta),\qquad
 \frac{d\phi}{d\theta}
 =\frac{a_Sb_S}{a_S^2\cos^2\theta+b_S^2\sin^2\theta}>0.
 \label{iv:ellipse:fases}
\end{equation}
Dans un système de coordonnées où les deux tangentes sont définies,
\(\tan\phi=(b_S/a_S)\tan\theta\). L’aire orientée balayée
à partir de la phase zéro est
\begin{equation}
 \mathcal A(\theta)=
 \frac12\int_0^\theta(QP'-PQ')\,dt=\hbar\theta,
 \qquad
 \mathcal A(1)=\hbar,\quad\mathcal A(2\pi)=h_s.
 \label{iv:ellipse:area-fase}
\end{equation}
\end{proposition}
\begin{proof}
La dérivée de l’argument d’un vecteur non nul est
\((QP'-PQ')/(Q^2+P^2)\). Pour la courbe
\eqref{iv:ellipse:curva}, le numérateur vaut
\(a_Sb_S=2\hbar\) ; cela démontre la dérivée et sa
positivité. Le quotient \(P/Q\) donne la relation entre tangentes ;
le relèvement continu conserve les quadrants et les tours.
L’intégrale vaut \(a_Sb_S\theta/2=\hbar\theta\).
Enfin, \(h_s=2\pi\hbar_s\) donne l’identité du
tour complet.

\end{proof}

Un radian de phase paramétrique balaie \(\hbar\). L’angle
polaire conserve sa propre dérivée dans
\eqref{iv:ellipse:fases}. En coordonnées polaires, la loi
équivalente est
\[
 d\mathcal A=\frac12\rho_{\partial}(\phi)^2\,d\phi,\qquad
 \rho_{\partial}(\phi)=
 \left(\frac{\cos^2\phi}{a_S^2}
       +\frac{\sin^2\phi}{b_S^2}\right)^{-1/2}.
\]
La formule radiale s’obtient en substituant
\((Q,P)=\rho(\cos\phi,\sin\phi)\) dans
\eqref{iv:ellipse:curva}. Avec l’orientation de
\(\gamma\), l’intégrale de la forme canonique de degré un conserve
le signe
\begin{equation}
 \oint_\gamma P\,dQ=-\pi a_Sb_S=-h_s.
 \label{iv:ellipse:accion-orientada}
\end{equation}
En effet, \(P\,dQ=-a_Sb_S\sin^2\theta\,d\theta\) et
l’intégrale de \(\sin^2\theta\) sur un tour est \(\pi\).

\subsection{Réciprocité de la forme et orientation symplectique}
\label{iv:ellipse:reciprocidad}

Nous définissons les matrices
\begin{equation}
 D(\eta)=\begin{pmatrix}e^\eta&0\\0&e^{-\eta}\end{pmatrix},
 \qquad
 S_2=\begin{pmatrix}0&1\\1&0\end{pmatrix},
 \qquad
 J_2=\begin{pmatrix}0&-1\\1&0\end{pmatrix}.
 \label{iv:ellipse:transportes}
\end{equation}

\begin{proposition}[Forme quadratique et action sous réciprocité]
\label{iv:ellipse:accion-reciproca}
Les deux transports satisfont
\(T^{\mathsf T}D(-\eta)T=D(\eta)\).
Pour \(\omega=dQ\wedge dP\), l’échange \(S_2\) est
antisymplectique et la rotation \(J_2\) est symplectique :
\[
 S_2^*\omega=-\omega,\qquad J_2^*\omega=\omega.
\]
Si \(\mathcal I(\gamma)=\oint_\gamma P\,dQ\) sur une
courbe fermée, alors
\begin{equation}
 \mathcal I(S_2\gamma)=-\mathcal I(\gamma),\qquad
 \mathcal I(J_2\gamma)=\mathcal I(\gamma).
 \label{iv:ellipse:signos-accion}
\end{equation}
\end{proposition}
\begin{proof}
La multiplication matricielle échange les deux entrées
diagonales de \(D\), ce qui démontre les deux congruences.
Les déterminants de \(S_2,J_2\) sont \(-1,1\) ;
on en obtient la transformation de la forme d’aire.
Pour \(\lambda=P\,dQ\), le calcul direct donne

\[
 S_2^*\lambda=d(PQ)-\lambda,\qquad
 J_2^*\lambda=\lambda-d(PQ).
\]
L’intégrale de la différentielle exacte sur la courbe fermée
est nulle, et les deux identités d’action sont prouvées.
\end{proof}

L’orientation distingue les deux transports bien qu’ils publient
la même forme quadratique, même dans le cas circulaire.
Le marquage des axes et le signe de l’action appartiennent
donc aux données conservées lors du passage à la
représentation réticulaire.

\subsection{Rayon normalisé et matrice de la lecture réticulaire}
\label{iv:ellipse:radio}

\begin{proposition}[Détermination et récupérabilité du rayon de forme]
\label{iv:ellipse:radio-prop}
L’ellipse étiquetée détermine un unique rayon positif
\begin{equation}
 R_\star=\sqrt{\frac{b_S}{a_S}}
 =e^{-\eta_{\mathrm{el}}/2}
 =\left(\frac{A-C^*}{A+C^*}\right)^{1/4},
 \qquad \tau_{\mathrm{el}}=iR_\star^2.
 \label{iv:ellipse:rstar}
\end{equation}
Le rapport angulaire se récupère au moyen de

\begin{equation}
 \frac{C^*}{A}=\frac{1-R_\star^4}{1+R_\star^4},
 \qquad
 R_\star^4=
 \frac{499\alpha-\eta_{\mathrm{ret}}}
      {501\alpha+\eta_{\mathrm{ret}}}.
 \label{iv:ellipse:rstar-compuesto}
\end{equation}
L’échange des demi-axes transforme
\((\eta_{\mathrm{el}},R_\star,\tau_{\mathrm{el}})\) en
\((-\eta_{\mathrm{el}},R_\star^{-1},-1/\tau_{\mathrm{el}})\)
et conserve \(a_Sb_S=2\hbar\).

\end{proposition}
\begin{proof}
La division des demi-axes donne \(b_S/a_S=e^{-\eta_{\mathrm{el}}}\).
L’application \(R\mapsto R^2\) est bijective sur les réels
positifs, de sorte que la racine positive est unique.
Pour \(\xi=C^*/A\in(-1,1)\),
\[
 2\operatorname{artanh}\xi
 =\log\frac{1+\xi}{1-\xi};
\]
par conséquent, \(R_\star^4=(1-\xi)/(1+\xi)\).
Isoler \(\xi\) démontre la première identité de
\eqref{iv:ellipse:rstar-compuesto}. Substituer
\(A=1000\alpha\) et
\(C^*=2(\eta_{\mathrm{ret}}+\alpha)\) donne
\[
 \frac{A-C^*}{A+C^*}
 =\frac{998\alpha-2\eta_{\mathrm{ret}}}
        {1002\alpha+2\eta_{\mathrm{ret}}}
 =\frac{499\alpha-\eta_{\mathrm{ret}}}
        {501\alpha+\eta_{\mathrm{ret}}}.
\]
L’échange inverse le rapport des demi-axes et change le
signe de \(\eta_{\mathrm{el}}\). L’identité
\(-1/(iR_\star^2)=iR_\star^{-2}\) donne le transport
modulaire. Le produit des demi-axes reste fixe.

\end{proof}

Sur la branche positive, \(0<R_\star<1\) ; le système de coordonnées de forme
inversée exprime \(R_\star^{-1}>1\). Le rapport angulaire
est récupéré à partir du rayon, tandis que les angles absolus
et leur histoire restent dans l’état enrichi.
La forme circulaire correspond à \(C^*/A=0\), avec
\(R_\star=1\). Changer la section positive d’action,
en maintenant la paire angulaire, modifie l’échelle commune des
demi-axes et conserve leur rapport.

\begin{proposition}[Forme réticulaire déterminée par l’ellipse]
\label{iv:ellipse:matriz-prop}
La matrice des demi-axes au carré normalisés est
\begin{equation}
 D(\eta_{\mathrm{el}})
 =\frac1{2\hbar}
    \begin{pmatrix}a_S^2&0\\0&b_S^2\end{pmatrix}
 =\begin{pmatrix}R_\star^{-2}&0\\0&R_\star^2\end{pmatrix}.
 \label{iv:ellipse:matriz}
\end{equation}
Son évaluation sur une paire entière produit

\begin{equation}
 E_\star(m,w)=
 \begin{pmatrix}m&w\end{pmatrix}
 D(\eta_{\mathrm{el}})
 \begin{pmatrix}m\\w\end{pmatrix}
 =\frac{m^2}{R_\star^2}+w^2R_\star^2.
 \label{iv:ellipse:energia}
\end{equation}
L’équation de l’ellipse multipliée par \(2\hbar\)
utilise, en revanche, la matrice inverse \(D(-\eta_{\mathrm{el}})\).
\end{proposition}
\begin{proof}
Les carrés des demi-axes de
\eqref{iv:ellipse:semiejes}, divisés par \(2\hbar\),
sont \(e^{\eta_{\mathrm{el}}}\) et
\(e^{-\eta_{\mathrm{el}}}\).
Par \eqref{iv:ellipse:rstar}, ce sont
\(R_\star^{-2}\) et \(R_\star^2\).
La multiplication par le vecteur entier démontre
\eqref{iv:ellipse:energia}. Enfin,
\(2\hbar/a_S^2=e^{-\eta_{\mathrm{el}}}\) et
\(2\hbar/b_S^2=e^{\eta_{\mathrm{el}}}\), ce qui
identifie la matrice de l’équation elliptique.

\end{proof}

Le rayon \(R_\star\) est une représentation sans dimension de
la forme étiquetée. Le rayon polaire
\(\rho_{\partial}\) appartient à la droite de racine d’action.
La représentation par un rayon spatial requiert sa mise en coordonnées
dimensionnelles ultérieure ; celle-ci conserve la forme déjà déterminée
et l’indépendance de sa construction par rapport à une
longueur prescrite.
